\section{Chains}

\import{order/order.tex}

The equality alternative in the definition makes chains insensitive to
whether the governing order relation is presented strictly or reflexively.

\begin{definition}\label{chain_in_relation}
    $C$ is a chain in $A$ with respect to $R$ iff
        $C\subseteq A$ and for all $x,y\in C$ we have
            $x=y$ or $x\mathrel{R}y$ or $y\mathrel{R}x$.
\end{definition}

\begin{definition}\label{maximal_chain_in_relation}
    $C$ is a maximal chain in $A$ with respect to $R$ iff
        $C$ is a chain in $A$ with respect to $R$ and
        for all $D$ such that $C\subseteq D$ and
            $D$ is a chain in $A$ with respect to $R$ we have $D=C$.
\end{definition}

\begin{definition}\label{inclusion_chain}
    $C$ is an inclusion chain iff for all $X,Y\in C$ we have
        $X\subseteq Y$ or $Y\subseteq X$.
\end{definition}

\begin{proposition}\label{emptyset_chain}
    $\emptyset$ is a chain in $A$ with respect to $R$.
\end{proposition}
\begin{proof}
    Follows by \auto{chain_in_relation,emptyset_subseteq,emptyset}.
\end{proof}

\begin{proposition}\label{singleton_chain}
    Suppose $a\in A$.
    Then $\{a\}$ is a chain in $A$ with respect to $R$.
\end{proposition}
\begin{proof}
    Follows by \auto{chain_in_relation,singleton_subset_intro,singleton_iff}.
\end{proof}

\begin{proposition}\label{subset_of_chain_is_chain}
    Suppose $D\subseteq C$.
    Suppose $C$ is a chain in $A$ with respect to $R$.
    Then $D$ is a chain in $A$ with respect to $R$.
\end{proposition}
\begin{proof}
    Follows by \auto{chain_in_relation,subseteq,subseteq_transitive}.
\end{proof}

\begin{proposition}\label{chain_tostrictorder_iff}
    $C$ is a chain in $A$ with respect to $R$ iff
        $C$ is a chain in $A$ with respect to $\tostrictorder{R}$.
\end{proposition}
\begin{proof}
    Follows by \auto{chain_in_relation,tostrictorder_iff}.
\end{proof}

\begin{proposition}\label{chain_toorder_iff}
    $C$ is a chain in $A$ with respect to $R$ iff
        $C$ is a chain in $A$ with respect to $\toorder{A}{R}$.
\end{proposition}
\begin{proof}
    Follows by \auto{chain_in_relation,toorder_iff}.
\end{proof}

\begin{proposition}\label{maximal_chain_tostrictorder_iff}
    $C$ is a maximal chain in $A$ with respect to $R$ iff
        $C$ is a maximal chain in $A$ with respect to $\tostrictorder{R}$.
\end{proposition}
\begin{proof}
    Follows by \auto{maximal_chain_in_relation,chain_tostrictorder_iff}.
\end{proof}

\begin{proposition}\label{maximal_chain_toorder_iff}
    $C$ is a maximal chain in $A$ with respect to $R$ iff
        $C$ is a maximal chain in $A$ with respect to $\toorder{A}{R}$.
\end{proposition}
\begin{proof}
    Follows by \auto{maximal_chain_in_relation,chain_toorder_iff}.
\end{proof}

\begin{proposition}\label{inclusion_chain_as_relation_chain}
    $C$ is an inclusion chain iff
        $C$ is a chain in $\pow{\unions{C}}$ with respect to
            $\subseteqrel{\pow{\unions{C}}}$.
\end{proposition}
\begin{proof}
    $C\subseteq\pow{\unions{C}}$
        by \auto{subseteq,pow_iff,unions_iff}.
    Show that if $C$ is an inclusion chain, then $C$ is a chain in
        $\pow{\unions{C}}$ with respect to
        $\subseteqrel{\pow{\unions{C}}}$.
    \begin{subproof}
        Assume $C$ is an inclusion chain.
        For all $X,Y\in C$ we have $X=Y$ or
            $X\mathrel{\subseteqrel{\pow{\unions{C}}}}Y$ or
            $Y\mathrel{\subseteqrel{\pow{\unions{C}}}}X$
            by \auto{inclusion_chain,subseteqrel,subseteq,pair_eq_iff}.
        Follows by \obv{chain_in_relation}.
    \end{subproof}
    Show that if $C$ is a chain in $\pow{\unions{C}}$ with respect to
        $\subseteqrel{\pow{\unions{C}}}$, then $C$ is an inclusion chain.
    \begin{subproof}
        Assume $C$ is a chain in $\pow{\unions{C}}$ with respect to
            $\subseteqrel{\pow{\unions{C}}}$.
        For all $X,Y\in C$ we have $X\subseteq Y$ or $Y\subseteq X$
            by \auto{chain_in_relation,subseteqrel,pair_eq_iff,subseteq_refl}.
        Follows by \obv{inclusion_chain}.
    \end{subproof}
    Follows by \auto.
\end{proof}

\begin{proposition}\label{inclusion_chain_in_family_iff}
    $C$ is a chain in $P$ with respect to $\subseteqrel{P}$ iff
        $C\subseteq P$ and $C$ is an inclusion chain.
\end{proposition}
\begin{proof}
    Show that if $C$ is a chain in $P$ with respect to $\subseteqrel{P}$,
        then $C\subseteq P$ and $C$ is an inclusion chain.
    \begin{subproof}
        Assume $C$ is a chain in $P$ with respect to $\subseteqrel{P}$.
        $C\subseteq P$ by \obv{chain_in_relation}.
        For all $X,Y\in C$ we have $X\subseteq Y$ or $Y\subseteq X$
            by \auto{chain_in_relation,subseteqrel,pair_eq_iff,subseteq_refl}.
        Thus $C$ is an inclusion chain by \obv{inclusion_chain}.
    \end{subproof}
    Show that if $C\subseteq P$ and $C$ is an inclusion chain,
        then $C$ is a chain in $P$ with respect to $\subseteqrel{P}$.
    \begin{subproof}
        Assume $C\subseteq P$ and $C$ is an inclusion chain.
        For all $X,Y\in C$ we have
            $X=Y$ or $X\mathrel{\subseteqrel{P}}Y$ or
                $Y\mathrel{\subseteqrel{P}}X$
            by \auto{inclusion_chain,subseteqrel,pair_eq_iff,subseteq}.
        Follows by \obv{chain_in_relation}.
    \end{subproof}
    Follows by \auto.
\end{proof}
